\documentclass[11pt]{article}
\usepackage{mathblatt}
% Präambel des Setzers (werkzeuge/setzer-praeambel.tex, wortgleich); dazu \gteil (Muster T2U)
\makeatletter
\setlength{\mbpfbe}{0mm}% keine Punkte (pruefung.md § 5: nur beim Prüfstein)
\newlength{\gEin}\setlength{\gEin}{0pt}
\let\gSchrift\relax
\renewcommand{\pfkreuzzeile}[1]{\par\noindent\hangindent3.2em\hangafter1\hspace*{1.6em}$\square$\ #1\par}
\newcommand{\gkreuz}[1]{$\square$\ #1\hspace{2em}}
\newcommand{\gkreuzzeile}[2]{\par\noindent\hangindent2.2em\hangafter1\mbox{#1\,$\square$}\ #2\par}% Buchstabe, Option: eine je Zeile (neben Skizze)
% links, was man ansieht, rechts, was man tut; Spaltengrenze überall an derselben Stelle (gemeinsam.md „Satz“)
% v0.17: Fußeintrag „Nr: Wert“ – kurz in einer Box (nie getrennt), lang umbrechbar mit der Nummer am Anfang
% (vorher \mbox: lange Lösungen liefen über den Rand, AS40 138 pt); robust, weil \fusshilfe in die .aux schreibt
\NewDocumentCommand{\gfuss}{m m}{\sbox0{#1:~#2}\ifdim\wd0<0.45\textwidth\mbox{#1:~#2}\else\mbox{#1:}\nobreak\ #2\fi}
\newsavebox{\gbox}\newlength{\gLinks}
% v0.10 (Lehrer 13.10.): keine starre Halbierung – die linke Spalte ist so breit wie die Figur (+ 6 mm), die
% Antwort steht direkt daneben; ist die Figur breiter als 0,62 Zeile, steht die Antwort darunter
\newcommand{\gzwei}[2]{\par\smallskip\noindent\sbox{\gbox}{#1}%
  \setlength{\gLinks}{\dimexpr\wd\gbox+6mm\relax}%
  \ifdim\gLinks>0.62\linewidth
    \usebox{\gbox}\par\smallskip\noindent\begin{minipage}{\linewidth}#2\end{minipage}%
  \else
    \begin{minipage}[t]{\gLinks}\vspace{0pt}\usebox{\gbox}\end{minipage}%
    \begin{minipage}[t]{\dimexpr\linewidth-\gLinks\relax}\vspace{0pt}\raggedright #2\end{minipage}%
  \fi\par}
% v0.9 (12.10., pruefung.md § 5 „Marke mit B, Bedingungen rechts“): drittes Argument von pfaufg = Bedingungen
% der Prüfung (v0.10: \pfohnetr = TR durchgestrichen, „FOR“), klein grau im rechten Rand auf der Grundlinie der Nummer, Zeilen mit \\
\newcommand{\gbedrand}[2]{\ifblank{#2}{}{\rlap{\hspace*{\dimexpr#1+2mm\relax}%
  \parbox[t]{15mm}{\raggedright\scriptsize\leavevmode\color{mbgrau}\strut#2}}}}
% Teil eines Bündels (a, b …): Marke links im Rand, Bedingungen rechts im Rand, auf der Zeile des Teils
\newcommand{\gteilrand}[2]{\llap{\mbpfrandmarke{#1}\hspace*{\mbpfnr}}\gbedrand{\linewidth}{#2}}
% Teil eines Bündels: Buchstabe in eigener Spalte, Inhalt eingerückt (\gEin = Einzug, damit die Spaltengrenze von
% \gzwei an derselben Stelle bleibt)
\newenvironment{gbteil}[3]{\par\smallskip\noindent\gteilrand{#2}{#3}\makebox[\mbpfteil][l]{\textbf{#1}}%
  \begin{minipage}[t]{\dimexpr\linewidth-\mbpfteil\relax}\setlength{\gEin}{\mbpfteil}\raggedright\strut\ignorespaces}%
 {\end{minipage}\par\addvspace{6pt}}% v0.18 (PG80 Nr. 1): Abstand, damit die Linie von a) nicht an b) rutscht
% Grafik höchstens so breit wie die Spalte (Bündel nebeneinander, Koordinatensystem ohne Karo)
\newsavebox{\gpbox}
\newcommand{\gpassend}[1]{\sbox{\gpbox}{#1}\ifdim\wd\gpbox>\linewidth\typeout{GPASSEND: verkleinert \the\wd\gpbox\space auf \the\linewidth}\resizebox{\linewidth}{!}{\usebox{\gpbox}}\else\usebox{\gpbox}\fi}
% Skizzenfeld (v0.10, Lehrer 13.10.): Karo ab dem linken Rand über die Zeile, „Skizze“ klein grau oben links im Feld
% (wie \pfskizzenfeld der Vorlage); der rechte Rand bleibt den Bedingungen
\newcommand{\gskizze}[1]{\par\smallskip\noindent\begin{tikzpicture}
  \clip (0,0) rectangle ({\linewidth-1mm},{#1*10mm});
  \draw[black!16,line width=0.3pt,step=5mm] (0,0) grid ({\linewidth},{#1*10mm});
  \node[anchor=north west,inner sep=2pt,font=\scriptsize,text=mbgrau,fill=white] at (0,{#1*10mm}) {Skizze};
  \end{tikzpicture}\par}
\RenewDocumentEnvironment{pfaufg}{m m m}{%
  \begin{lrbox}{\mbaufgabebox}\begin{minipage}[t]{\linewidth}%
  \gSchrift\noindent\hspace*{\gEin}\mbpfrandmarke{#1}\gbedrand{\linewidth-\gEin}{#3}%
  \begin{minipage}[t]{\mbpfnr}\strut\textbf{#2}\end{minipage}%
  \begin{minipage}[t]{\dimexpr\linewidth-\gEin-\mbpfnr\relax}\raggedright\strut\ignorespaces}%
 {\end{minipage}%
  \end{minipage}\end{lrbox}%
  \setlength{\mbaufgabehoehe}{\dimexpr\ht\mbaufgabebox+\dp\mbaufgabebox\relax}%
  \ifdefempty{\gkopfnach}{\par\addvspace{5pt}\Needspace*{\mbaufgabehoehe}}{%
    % Stufenkopf, der mit dieser (ersten) Aufgabe zusammenbleibt: Platz für beide verlangen, dann den Kopf setzen
    % (12.10., Quadratische: ein Kopf stand allein unten, die Aufgabe mit Koordinatensystem auf der nächsten Seite)
    % v0.17 (AS40 „Terme deuten“ allein unten): hinter dem Kopf steht das Sprungziel (\special, ein Knoten, nach
    % dem die Seite brechen durfte) – darum \nobreak nach dem Kopf; Vorhalt 56 pt = Kopf mit Linie und Abständen
    \par\Needspace*{\dimexpr\mbaufgabehoehe+56pt\relax}\gkopfnach\par\nobreak\gdef\gkopfnach{}\addvspace{5pt}}%
  \noindent\usebox{\mbaufgabebox}\par\addvspace{5pt}}
% v0.18 (plan.md § 10, 15.10. l): Abschnittsbuchstabe A, B, C vor dem Namen im Kopf – als Vorsatz, damit der
% Name im TeX gleich bleibt (Bestellkarte, Sprungziele lesen ihn)
\def\gbstakt{}\newcommand{\gbstvor}[1]{\gdef\gbstakt{#1}}
\NewCommandCopy\gpfstufekopfalt\pfstufekopf
\RenewDocumentCommand{\pfstufekopf}{o m m}{\IfValueTF{#1}{\gpfstufekopfalt[#1]}{\gpfstufekopfalt}%
  {\ifx\gbstakt\empty\else\gbstakt\enspace\fi #2}{#3}\gdef\gbstakt{}}
% v0.18 (16.10. e): Begriff grau, ohne Nummer, ohne Rahmen über der ersten Aufgabe der Stufe
\newcommand{\gbegriff}[1]{\par\nobreak\noindent\hspace*{\mbpfnr}\parbox[t]{\dimexpr\linewidth-\mbpfnr\relax}{\raggedright\small\color{mbgrau}#1}\par\nobreak\addvspace{3pt}}
% v0.18 (16.10. g): Lösungsblatt – Kopf mit kräftigem Strich, Abschnitt klein grau als Zwischenzeile
\newcommand{\glkopf}[1]{\noindent{\large\bfseries Lösungen \textperiodcentered{} #1}\par\vspace{3pt}\hrule height 1.2pt\par\medskip}
\newcommand{\glzwischen}[1]{\par\addvspace{8pt}\Needspace*{4\baselineskip}\noindent{\footnotesize\color{mbgrau}#1}\par\nobreak\vspace{1pt}}
% v0.18 (15.10. o, w): Blattende – Kern ohne Rahmen, dann graue Zeile „Mehr: n Aufgaben“
\newcommand{\gkern}[1]{\par\addvspace{14pt}\Needspace*{6\baselineskip}\noindent{\bfseries Kern}\par\nobreak\smallskip
  {\small\raggedright #1\par}}
\newcommand{\gmehr}[1]{\par\addvspace{10pt}\noindent{\footnotesize\color{mbgrau}Mehr: #1 Aufgaben}\par}
% Stufenkopf vormerken: \gkopfvor{<Kopf>} – die nächste pfaufg setzt ihn, sobald beide zusammen Platz haben
\def\gkopfnach{}
\newcommand{\gkopfvor}[1]{\gdef\gkopfnach{#1}}
% Teilaufgabe mit Skizze links in gzwei: Buchstabe in eigener Spalte, darunter Text und Skizze
\newcommand{\gteil}[2]{\makebox[\mbpfteil][l]{\textbf{#1}}\parbox[t]{\dimexpr0.5\textwidth-\gEin-\mbpfnr-5mm-\mbpfteil\relax}{\raggedright #2}}
% Stufenkopf der Vorlage (mathblatt.sty 2026-10-12: Linie, Abstand, 13 pt, \pfjahre) ohne ihren Vorhalt von
% 14 Zeilen: der Setzer fordert vor dem Kopf 0,16 Seitenhöhe (Kopf + erste Aufgabe), mehr nicht
\setlength{\mbpfstufevorhalt}{0pt}
\makeatother
\newcommand{\karo}[1]{\par\smallskip\noindent\begin{tikzpicture}
  \clip (0,0) rectangle ({\linewidth-1mm},{#1*10mm});
  \draw[black!16,line width=0.3pt,step=5mm] (0,0) grid ({\linewidth},{#1*10mm});
  \end{tikzpicture}\par}
\newcommand{\kk}[1]{$\square$\,#1\hspace{0.9em}}
\newcommand{\linie}[1]{\,{\color{black!45}\rule[-1pt]{#1}{0.4pt}}\,}
\newcommand{\pkt}{\textperiodcentered{}}
% v0.14 (14.10.): Stufenkopf mit fünf Punkten (Option = Jahre für die Bestellkarte) oder „zuletzt 19“
\NewDocumentCommand{\gpunkte}{o m}{#2}
\newcommand{\gzuletzt}[1]{zuletzt #1}
% v0.7 (12.10.): Schülerfrage unter dem Stufenkopf; Ankreuzoption als Bild (Buchstabe und Kästchen darunter)
\newcommand{\gfrage}[1]{\par\nobreak\noindent\hspace*{\mbpfnr}{\small\itshape #1}\par\nobreak\addvspace{2pt}}
\newlength{\gbildbreite}\setlength{\gbildbreite}{0.32\linewidth}% Setzer setzt sie je Reihe (vier Bilder: 0,24)
\newcommand{\gbildoption}[3]{\begin{minipage}[t][#3][b]{\gbildbreite}\centering #2\par\smallskip\mbox{#1\,$\square$}\end{minipage}}
% Lineare (12.10.): Bild mit Buchstabe ohne Kästchen (Diagramm zuordnen); Aussage mit „wahr / falsch“ zum Ankreuzen
\newcommand{\gbildname}[3]{\begin{minipage}[t][#3][b]{\gbildbreite}\centering #2\par\smallskip\mbox{\textbf{#1}}\end{minipage}}
\newcommand{\gwahrfalsch}[2]{\par\noindent\hangindent2.2em\hangafter1\makebox[2.2em][l]{(#1)}#2\hfill\mbox{wahr $\square$\quad falsch $\square$}\par}
\usepackage{graphicx}\graphicspath{{/home/claude/aufgabenbank-privat/}}
\begin{document}

\pfheftstil
\pflaufkopf[P10 \pkt{} Prüfung]{Lineare Funktionen}{PA10-8}\pfseitevon
\pfheftkopf{Lineare Funktionen}{}
\pfstufenzeile{A~Graph zeichnen \pkt{} B~Gleichung aus Graph \pkt{} C~Gleichung aus zwei Punkten \pkt{} D~Nullstelle gesucht \pkt{} E~Punktprobe \pkt{} F~Gleichung aus Text}
% Formelkasten (15.10. t, 16.10. a)
{\small% duden-s29-y-f-x-m-2 – y = f(x) = m· x + n für m ≠ 0 (Duden S. 29), aus duden-sek1.md
\providecommand\fsf[1]{\leavevmode\setbox0\hbox{\everymath{\displaystyle}#1}\ifdim\wd0>\linewidth\setbox0\hbox{\resizebox{\linewidth}{!}{\box0}}\fi
 \vrule height\dimexpr\ht0+1.5pt\relax depth\dimexpr\dp0+1.5pt\relax width0pt\box0}
\providecommand\fstab[1]{\leavevmode\setbox0\hbox{#1}\ifdim\wd0>\linewidth\resizebox{\linewidth}{!}{\box0}\else\box0\fi}
\providecommand\fssic[1]{{\color{red!65!black}\footnotesize[sic?#1]}}
\providecommand\fsunl[1]{{\color{red!65!black}\footnotesize[unleserlich#1]}}
\providecommand\fsfigur[1]{{\footnotesize\color{mbgrau}\itshape Figur im Original: #1\par}}
\providecommand\fsbild[1]{\leavevmode\hfill#1\par}
\providecommand\fsbildbreite{4.6cm}
\providecommand\fszeilebild[6]{\par\noindent\begin{minipage}[t]{\dimexpr\fsbreite-#6-6pt}\vspace{0pt}\renewcommand\fsbreite{\linewidth}\fszeile{#1}{#2}{#3}{}\end{minipage}\hfill
 \begin{minipage}[t]{#6}\vspace{3pt}\includegraphics[width=#6]{#5}\end{minipage}\par}
\usetikzlibrary{angles}\tikzset{fsfig/.style={line width=0.7pt,baseline=(current bounding box.north),font=\small},
 fswinkel/.style={draw=mbgrau,line width=0.5pt,angle radius=5.5mm,angle eccentricity=0.62,pic text options={font=\footnotesize}},
 fsrecht/.style={draw=mbgrau,line width=0.5pt,angle radius=2.6mm,angle eccentricity=0.5,pic text={\tiny$\bullet$}},
 fslinie/.style={font=\scriptsize,sloped,inner sep=1.5pt},fsflaeche/.style={fill=mbgrau,fill opacity=0.18,draw=none}}
\providecommand\fspkt[1]{\fill #1 circle (1.4pt);}
\providecommand\fshinweis[1]{{\footnotesize\color{mbgrau}#1\par}}
\providecommand\fspunkt{\ensuremath{\scriptstyle\blacksquare}\,}
\providecommand\fsfuenf{\ensuremath{\rlap{$||||$}\diagup}}
\providecommand\fsquelle[1]{\hfill{\footnotesize\color{mbgrau}#1}}
\providecommand\fsbreite{\linewidth}
\providecommand\fszeile[4]{\par\noindent{\renewcommand\arraystretch{1.25}\setlength\lineskip{4pt}\setlength\lineskiplimit{3pt}%
 \ifblank{#2}%
 {\begin{tabular}{>{\columncolor{mbkasten}\raggedright\hspace{0pt}}p{\dimexpr0.25\fsbreite-2\tabcolsep}>{\raggedright\arraybackslash}p{\dimexpr0.75\fsbreite-2\tabcolsep}}%
  #1 & #3\ifblank{#4}{}{\\ & #4}\end{tabular}}%
 {\begin{tabular}{>{\columncolor{mbkasten}\raggedright\hspace{0pt}}p{\dimexpr0.25\fsbreite-2\tabcolsep}>{\raggedright}p{\dimexpr0.33\fsbreite-2\tabcolsep}>{\raggedright\arraybackslash}p{\dimexpr0.42\fsbreite-2\tabcolsep}}%
  #1 & #2 & #3\ifblank{#4}{}{\\ & \multicolumn{2}{>{\raggedright\arraybackslash}p{\dimexpr0.75\fsbreite-2\tabcolsep}}{#4}}\end{tabular}}}\par}
\providecommand\fskastentitel[1]{\par\noindent\textbf{#1}\par}
\uebersichtskasten{\renewcommand\fsbreite{\linewidth}
\fszeilebild{$y = f(x) = m\cdot x + n$ für $m \neq 0$}{Gerade durch P(0; n)\newline Definitionsbereich: $x \in \mathbb{R}$\newline Wertebereich: $y \in \mathbb{R}$\newline n: Abschnitt auf y-Achse\newline Nullstelle: $x_0 = -\frac{n}{m}$ (blau)}{}{}{formelsammlung/figuren/duden-bild/duden-s29-y-f-x-m-2.png}{3.89cm}
}
}\vspace{-4pt}\par\noindent{\footnotesize\color{mbgrau}Formeln und Werte S.~29}\par\medskip

% ===================== Stufe: Gerade zeichnen =====================
\gkopfvor{\gbstvor{A}\pfstufekopf{Graph zeichnen}{\gpunkte[26 25 24 23 22]{\textcolor{black!80}{$\bullet$}\,\textcolor{black!80}{$\bullet$}\,\textcolor{black!80}{$\bullet$}\,\textcolor{black!80}{$\bullet$}\,\textcolor{black!80}{$\bullet$}}}\special{pdf:dest (stufe-graph-zeichnen) [@thispage /XYZ @xpos @ypos null]}}
% kein Einstieg: keine Schülerfrage in der Gliederung (Schülerfrage, -Fälle, -Form, -Weg, -Skizze; oder Frage) – Einstieg nur von dort, nie aus der Bank (pruefung.md § 5, 12.10.)
% 2025-OS-B1f-m1 (Original 2025-OS-B1f, Schwere I, rang 1; herausgelöste Fassung: keine)
\begin{pfaufg}{’25~\pkt{}~1f}{1.}{Basis}
Kreuze das Bild an, das den Graphen einer linearen Funktion zeigt.
\par\smallskip\noindent\setlength{\gbildbreite}{\dimexpr0.96\linewidth/4\relax}\sbox{\gbox}{\vbox{\begin{tikzpicture}[x=1.1cm,y=1.1cm,line width=0.6pt]\draw[-{Stealth[length=1.5mm]}] (-0.8,0)--(0.95,0);\draw[-{Stealth[length=1.5mm]}] (0,-0.6)--(0,1.05);\draw[line width=0.8pt] (-0.8,0.94)--(0.9,-0.42);\end{tikzpicture}\par\begin{tikzpicture}[x=1.1cm,y=1.1cm,line width=0.6pt]\draw[-{Stealth[length=1.5mm]}] (-0.8,0)--(0.95,0);\draw[-{Stealth[length=1.5mm]}] (0,-0.6)--(0,1.05);\draw[line width=0.8pt,domain=0.1:0.9,samples=40,smooth] plot (\x,{0.09/\x});\end{tikzpicture}\par\begin{tikzpicture}[x=1.1cm,y=1.1cm,line width=0.6pt]\draw[-{Stealth[length=1.5mm]}] (-0.8,0)--(0.95,0);\draw[-{Stealth[length=1.5mm]}] (0,-0.6)--(0,1.05);\draw[line width=0.8pt,domain=-0.6:0.6,samples=40,smooth] plot (\x,{3*\x*\x-0.4});\end{tikzpicture}\par\begin{tikzpicture}[x=1.1cm,y=1.1cm,line width=0.6pt]\draw[-{Stealth[length=1.5mm]}] (-0.8,0)--(0.95,0);\draw[-{Stealth[length=1.5mm]}] (0,-0.6)--(0,1.05);\draw[line width=0.8pt] (0.1,-0.3)--(0.5,0.5)--(0.1,0.7);\end{tikzpicture}}}\gbildoption{}{\begin{tikzpicture}[x=1.1cm,y=1.1cm,line width=0.6pt]\draw[-{Stealth[length=1.5mm]}] (-0.8,0)--(0.95,0);\draw[-{Stealth[length=1.5mm]}] (0,-0.6)--(0,1.05);\draw[line width=0.8pt] (-0.8,0.94)--(0.9,-0.42);\end{tikzpicture}}{\dimexpr\ht\gbox/4+\baselineskip\relax}\hfill\gbildoption{}{\begin{tikzpicture}[x=1.1cm,y=1.1cm,line width=0.6pt]\draw[-{Stealth[length=1.5mm]}] (-0.8,0)--(0.95,0);\draw[-{Stealth[length=1.5mm]}] (0,-0.6)--(0,1.05);\draw[line width=0.8pt,domain=0.1:0.9,samples=40,smooth] plot (\x,{0.09/\x});\end{tikzpicture}}{\dimexpr\ht\gbox/4+\baselineskip\relax}\hfill\gbildoption{}{\begin{tikzpicture}[x=1.1cm,y=1.1cm,line width=0.6pt]\draw[-{Stealth[length=1.5mm]}] (-0.8,0)--(0.95,0);\draw[-{Stealth[length=1.5mm]}] (0,-0.6)--(0,1.05);\draw[line width=0.8pt,domain=-0.6:0.6,samples=40,smooth] plot (\x,{3*\x*\x-0.4});\end{tikzpicture}}{\dimexpr\ht\gbox/4+\baselineskip\relax}\hfill\gbildoption{}{\begin{tikzpicture}[x=1.1cm,y=1.1cm,line width=0.6pt]\draw[-{Stealth[length=1.5mm]}] (-0.8,0)--(0.95,0);\draw[-{Stealth[length=1.5mm]}] (0,-0.6)--(0,1.05);\draw[line width=0.8pt] (0.1,-0.3)--(0.5,0.5)--(0.1,0.7);\end{tikzpicture}}{\dimexpr\ht\gbox/4+\baselineskip\relax}\par
\fusshilfe{\gfuss{1}{die einzige Gerade; B Hyperbelast, C Parabel, D kein Funktionsgraph}}
\end{pfaufg}
% 2024-OS-B1i-m1 (Original 2024-OS-B1i, Schwere I, rang 2; herausgelöste Fassung: keine)
\begin{pfaufg}{’24~\pkt{}~1i}{2.}{Basis}
Gegeben ist die Gerade $y = 1{,}5x + 1$. Zeichne den Punkt ein, an dem sie die y-Achse trifft.
\par\smallskip\noindent\begin{ksys}[xmin=-3,xmax=3,ymin=-3,ymax=3,karo=0.5]\end{ksys}\par
\fusshilfe{\gfuss{2}{Punkt (0|1)}}
\end{pfaufg}
% 2024-OS-K3a-m1 (Original 2024-OS-K3a, Schwere I, rang 3; herausgelöste Fassung: keine)
\begin{pfaufg}{nach~P10\\’24~\pkt{}~3a}{3.}{}
Die lineare Funktion $f$ hat die Gleichung $f(x) = 2x + 3$. Zeichne ihren Graphen in das Koordinatensystem.
\par\smallskip\noindent\gpassend{\begin{ksys}[xmin=-4,xmax=4,ymin=-5,ymax=8,karo=0.5]\end{ksys}}\par
\fusshilfe{\gfuss{3}{Gerade durch (0|3) und (1|5)}}
\end{pfaufg}

% ===================== Stufe: Gleichung aus Graph =====================
\gkopfvor{\gbstvor{B}\pfstufekopf{Gleichung aus Graph}{\gzuletzt{19}}\special{pdf:dest (stufe-gleichung-aus-graph) [@thispage /XYZ @xpos @ypos null]}}
% kein Einstieg: keine Schülerfrage in der Gliederung (Schülerfrage, -Fälle, -Form, -Weg, -Skizze; oder Frage) – Einstieg nur von dort, nie aus der Bank (pruefung.md § 5, 12.10.)
% Teile eines Originals (2019-OS-K2): 2019-OS-K2b-m1, 2019-OS-K2c-m1
\begin{pfaufg}{nach~P10\\’19~\pkt{}~2}{4.}{}
In einem Koordinatensystem sind die Graphen $f$, $g$ und $h$ von drei linearen Funktionen gezeichnet.
% 2019-OS-K2b-m1 (Original 2019-OS-K2b, Schwere I, rang 11; herausgelöste Fassung: keine)
\begin{gbteil}{a)}{}{}
Schreib zu jeder Eigenschaft den Graphen ($f$, $g$ oder $h$), der sie hat: „Der Graph hat die Steigung $-3$.“ „Der Graph liegt parallel zur x-Achse.“
\gzwei{\begin{ksys}[xmin=-5,xmax=5,ymin=-5,ymax=5,karo=0.5]\gerade[0.6]{-3}{3}{f}\gerade[1]{1}{3}{g}\gerade[3.5]{0}{-1}{h}\end{ksys}}{\vspace{1mm}\par\smallskip\par\noindent \hfill Steigung $-$3: \linie{30mm}\hphantom{\,cm²}\par\noindent \hfill parallel zur x-Achse: \linie{30mm}\hphantom{\,cm²}}
\end{gbteil}
\end{pfaufg}
\begin{pfaufg}{}{}{}
% 2019-OS-K2c-m1 (Original 2019-OS-K2c, Schwere II, rang 13; herausgelöste Fassung: keine)
\begin{gbteil}{b)}{}{}
Ordne jedem Graphen $f$, $g$ und $h$ eine Gleichung aus dieser Liste zu: $y = -3x + 3$, $y = -\tfrac{1}{3}x + 3$, $y = -1$, $y = x + 3$, $y = -3x$, $y = x - 3$.
\gzwei{\begin{ksys}[xmin=-5,xmax=5,ymin=-5,ymax=5,karo=0.5]\gerade[0.6]{-3}{3}{f}\gerade[1]{1}{3}{g}\gerade[3.5]{0}{-1}{h}\end{ksys}}{\vspace{1mm}\par\smallskip\par\noindent \hfill f: \linie{30mm}\hphantom{\,cm²}\par\noindent \hfill g: \linie{30mm}\hphantom{\,cm²}\par\noindent \hfill h: \linie{30mm}\hphantom{\,cm²}}
\end{gbteil}
\fusshilfe{\gfuss{4}{a)~h; b)~$-$1}}
\end{pfaufg}

% ===================== Stufe: Gleichung aus zwei Punkten =====================
\gkopfvor{\gbstvor{C}\pfstufekopf{Gleichung aus zwei Punkten}{\gpunkte[25]{\textcolor{black!25}{$\circ$}\,\textcolor{black!80}{$\bullet$}\,\textcolor{black!25}{$\circ$}\,\textcolor{black!25}{$\circ$}\,\textcolor{black!25}{$\circ$}}}\special{pdf:dest (stufe-gleichung-aus-zwei-punkten) [@thispage /XYZ @xpos @ypos null]}}
% kein Einstieg: keine Schülerfrage in der Gliederung (Schülerfrage, -Fälle, -Form, -Weg, -Skizze; oder Frage) – Einstieg nur von dort, nie aus der Bank (pruefung.md § 5, 12.10.)
% 2025-OS-K5a-m1 (Original 2025-OS-K5a, Schwere II, rang 6; herausgelöste Fassung: keine)
\begin{pfaufg}{nach~P10\\’25~\pkt{}~5a}{5.}{}
Der Graph einer linearen Funktion $f$ geht durch die Punkte $A(-2\,|\,5)$ und $B(3\,|\,-2{,}5)$. Zeichne den Graphen von $f$ in das Koordinatensystem. Entscheide, ob die Aussagen wahr oder falsch sind:
\gzwei{\begin{ksys}[xmin=-4,xmax=4,ymin=-4,ymax=8,karo=0.5]\end{ksys}}{\vspace{1mm}\gwahrfalsch{1}{Der Graph von $f$ steigt monoton.}\gwahrfalsch{2}{Der Graph von $f$ schneidet die y-Achse in $(0\,|\,2)$}\par\smallskip\noindent Gib eine Gleichung von $f$ an.\par\smallskip\noindent y $=$\linie{40mm}}
\fusshilfe{\gfuss{5}{Gerade durch A und B; Aussage 1 falsch, Aussage 2 wahr; $-$1,5x + 2}}
\end{pfaufg}

% ===================== Stufe: Nullstelle =====================
\gkopfvor{\gbstvor{D}\pfstufekopf{Nullstelle gesucht}{\gpunkte[22]{\textcolor{black!25}{$\circ$}\,\textcolor{black!25}{$\circ$}\,\textcolor{black!25}{$\circ$}\,\textcolor{black!25}{$\circ$}\,\textcolor{black!80}{$\bullet$}}}\special{pdf:dest (stufe-nullstelle-gesucht) [@thispage /XYZ @xpos @ypos null]}}
% kein Einstieg: keine Schülerfrage in der Gliederung (Schülerfrage, -Fälle, -Form, -Weg, -Skizze; oder Frage) – Einstieg nur von dort, nie aus der Bank (pruefung.md § 5, 12.10.)
% 2022-OS-K3a-m1 (Original 2022-OS-K3a, Schwere I, rang 7; herausgelöste Fassung: keine)
\begin{pfaufg}{nach~P10\\’22~\pkt{}~3a}{6.}{}
Gegeben ist die lineare Funktion $f(x) = -2x + 5$. Zeichne den Graphen von $f$ in das Koordinatensystem. Gib die Nullstelle von $f$ an.
\gzwei{\begin{ksys}[xmin=-4,xmax=5,ymin=-4,ymax=7,karo=0.5]\end{ksys}}{\vspace{1mm}\par\smallskip\noindent \mbox{x $=$\linie{30mm}}}
\fusshilfe{\gfuss{6}{Gerade durch (0|5) und (1|3); 2,5}}
\end{pfaufg}

% ===================== Stufe: Punktprobe =====================
\gkopfvor{\gbstvor{E}\pfstufekopf{Punktprobe}{\gpunkte[26 25 24 23 22]{\textcolor{black!80}{$\bullet$}\,\textcolor{black!80}{$\bullet$}\,\textcolor{black!80}{$\bullet$}\,\textcolor{black!80}{$\bullet$}\,\textcolor{black!80}{$\bullet$}}}\special{pdf:dest (stufe-punktprobe) [@thispage /XYZ @xpos @ypos null]}}
% kein Einstieg: keine Schülerfrage in der Gliederung (Schülerfrage, -Fälle, -Form, -Weg, -Skizze; oder Frage) – Einstieg nur von dort, nie aus der Bank (pruefung.md § 5, 12.10.)
% 2023-OS-B1i-m1 (Original 2023-OS-B1i, Schwere I, rang 8; herausgelöste Fassung: keine)
\begin{pfaufg}{’23~\pkt{}~1i}{7.}{Basis}
Die Gerade $f$ hat die Gleichung $f(x) = -5x + 2$. Einer dieser Punkte liegt nicht auf $f$. Kreuze ihn an:
\par\smallskip\noindent\pfkreuz{}{$(-2\,|\,12)$}\pfkreuz{}{$(3\,|\,-12)$}\pfkreuz{}{$(4\,|\,-18)$.}
\fusshilfe{\gfuss{7}{$(3\,|\,-12)$}}
\end{pfaufg}
% 2026-FOR-K5b-m1 (Original 2026-FOR-K5b, Schwere I, rang 10; herausgelöste Fassung: keine)
\begin{pfaufg}{nach~P10\\’26~\pkt{}~5b}{8.}{}
Prüfe mit einer Rechnung, ob der Punkt $P(-2\,|\,7)$ auf dem Graphen von $f(x) = -3x + 1$ liegt.
\karo{2}
\fusshilfe{\gfuss{8}{P liegt auf dem Graphen}}
\end{pfaufg}

% ===================== Stufe: Gleichung aus Text =====================
\Needspace*{0.16\textheight}\gbstvor{F}\pfstufekopf{Gleichung aus Text}{\gpunkte[23 22]{\textcolor{black!25}{$\circ$}\,\textcolor{black!25}{$\circ$}\,\textcolor{black!25}{$\circ$}\,\textcolor{black!80}{$\bullet$}\,\textcolor{black!80}{$\bullet$}}}\special{pdf:dest (stufe-gleichung-aus-text) [@thispage /XYZ @xpos @ypos null]}
% Schülerfrage der Stufe (Gliederung; Einstieg nur von dort, nie aus der Bank – 12.10.)
\begin{pfaufg}{}{9.}{}
Welche Zahl ist m, welche n? Rechne nicht.
\par\smallskip
\par\noindent\makebox[\mbpfteil][l]{\textbf{a)}}Handytarif: 8 € im Monat und 0,09 € je Minute\hfill\linie{30mm}
\par\noindent\makebox[\mbpfteil][l]{\textbf{b)}}Kerze: 25 cm hoch, brennt je Stunde 2 cm ab\hfill\linie{30mm}
\par\fusshilfe{\gfuss{9}{a) 8 ist n, 0,09 ist m; b) 25 ist n, $-$2 ist m}}
\end{pfaufg}
% Teile eines Originals (2022-OS-K6): 2022-OS-K6a-m1, 2022-OS-K6b-m1
\begin{pfaufg}{nach~P10\\’22~\pkt{}~6}{10.}{}
Maxi und Paula haben zum 16. Geburtstag Geld bekommen und legen jeden Monat etwas dazu, um den Führerschein zu bezahlen. Maxi beginnt mit 1 360 € und spart 24 € im Monat. Paula beginnt mit 820 € und spart 65 € im Monat.
% 2022-OS-K6a-m1 (Original 2022-OS-K6a, Schwere I, rang 12; herausgelöste Fassung: keine)
\begin{gbteil}{a)}{}{}
Nach einem Jahr hat Paula 1 600 € gespart. Berechne, wie viel Geld Maxi nach einem Jahr hat.
\karo{2}
\end{gbteil}
\end{pfaufg}
\begin{pfaufg}{}{}{}
% 2022-OS-K6b-m1 (Original 2022-OS-K6b, Schwere II, rang 13; herausgelöste Fassung: keine)
\begin{gbteil}{b)}{}{FOR}
Stelle eine Gleichung für Paulas Guthaben auf ($y$: Guthaben in €, $x$: Zahl der Monate). Für Fahrschule und Prüfung braucht Paula 1 800 €. Bestimme, wie viele Monate sie mindestens sparen muss.
\karo{3}
\end{gbteil}
\fusshilfe{\gfuss{10}{a)~1648 €; b)~65x + 820; 16 Monate}}
\end{pfaufg}
% Teile eines Originals (2023-OS-K3): 2023-OS-K3a-m1, 2023-OS-K3b-m1
\begin{pfaufg}{nach~P10\\’23~\pkt{}~3}{11.}{}
Herr Mert braucht für seinen Umzug einen Transporter. Angebot 1: 25 € Miete je Tag, 320 km insgesamt frei, jeder weitere Kilometer kostet 0,35 €. Angebot 2: einmalig 115 € Miete für eine Woche, jeder Kilometer kostet 0,08 €.
% 2023-OS-K3a-m1 (Original 2023-OS-K3a, Schwere II, rang 14; herausgelöste Fassung: keine)
\begin{gbteil}{a)}{}{}
Die Kosten für einen Miettag sollen für beide Angebote in einem Diagramm dargestellt werden. Gib zu jedem Angebot das passende Diagramm an.
\par\smallskip\noindent\setlength{\gbildbreite}{\dimexpr0.96\linewidth/4\relax}\sbox{\gbox}{\vbox{\begin{tikzpicture}[x=1.3cm,y=1.3cm,line width=0.6pt]\draw[-{Stealth[length=1.5mm]}] (0,0)--(1.15,0) node[below left,font=\tiny] {km};\draw[-{Stealth[length=1.5mm]}] (0,0)--(0,1.05) node[below right,font=\tiny] {€};\draw[line width=0.8pt] (0,0)--(1,0.85);\end{tikzpicture}\par\begin{tikzpicture}[x=1.3cm,y=1.3cm,line width=0.6pt]\draw[-{Stealth[length=1.5mm]}] (0,0)--(1.15,0) node[below left,font=\tiny] {km};\draw[-{Stealth[length=1.5mm]}] (0,0)--(0,1.05) node[below right,font=\tiny] {€};\draw[line width=0.8pt] (0,0.35)--(1,0.6);\end{tikzpicture}\par\begin{tikzpicture}[x=1.3cm,y=1.3cm,line width=0.6pt]\draw[-{Stealth[length=1.5mm]}] (0,0)--(1.15,0) node[below left,font=\tiny] {km};\draw[-{Stealth[length=1.5mm]}] (0,0)--(0,1.05) node[below right,font=\tiny] {€};\draw[line width=0.8pt] (0,0.3)--(0.45,0.3)--(1,0.85);\end{tikzpicture}\par\begin{tikzpicture}[x=1.3cm,y=1.3cm,line width=0.6pt]\draw[-{Stealth[length=1.5mm]}] (0,0)--(1.15,0) node[below left,font=\tiny] {km};\draw[-{Stealth[length=1.5mm]}] (0,0)--(0,1.05) node[below right,font=\tiny] {€};\draw[line width=0.8pt] (0,0.7)--(0.45,0.7)--(1,0.25);\end{tikzpicture}}}\gbildname{A}{\begin{tikzpicture}[x=1.3cm,y=1.3cm,line width=0.6pt]\draw[-{Stealth[length=1.5mm]}] (0,0)--(1.15,0) node[below left,font=\tiny] {km};\draw[-{Stealth[length=1.5mm]}] (0,0)--(0,1.05) node[below right,font=\tiny] {€};\draw[line width=0.8pt] (0,0)--(1,0.85);\end{tikzpicture}}{\dimexpr\ht\gbox/4+\baselineskip\relax}\hfill\gbildname{B}{\begin{tikzpicture}[x=1.3cm,y=1.3cm,line width=0.6pt]\draw[-{Stealth[length=1.5mm]}] (0,0)--(1.15,0) node[below left,font=\tiny] {km};\draw[-{Stealth[length=1.5mm]}] (0,0)--(0,1.05) node[below right,font=\tiny] {€};\draw[line width=0.8pt] (0,0.35)--(1,0.6);\end{tikzpicture}}{\dimexpr\ht\gbox/4+\baselineskip\relax}\hfill\gbildname{C}{\begin{tikzpicture}[x=1.3cm,y=1.3cm,line width=0.6pt]\draw[-{Stealth[length=1.5mm]}] (0,0)--(1.15,0) node[below left,font=\tiny] {km};\draw[-{Stealth[length=1.5mm]}] (0,0)--(0,1.05) node[below right,font=\tiny] {€};\draw[line width=0.8pt] (0,0.3)--(0.45,0.3)--(1,0.85);\end{tikzpicture}}{\dimexpr\ht\gbox/4+\baselineskip\relax}\hfill\gbildname{D}{\begin{tikzpicture}[x=1.3cm,y=1.3cm,line width=0.6pt]\draw[-{Stealth[length=1.5mm]}] (0,0)--(1.15,0) node[below left,font=\tiny] {km};\draw[-{Stealth[length=1.5mm]}] (0,0)--(0,1.05) node[below right,font=\tiny] {€};\draw[line width=0.8pt] (0,0.7)--(0.45,0.7)--(1,0.25);\end{tikzpicture}}{\dimexpr\ht\gbox/4+\baselineskip\relax}\par
\par\noindent\hfill Angebot 1 $=$\linie{30mm}
\par\noindent\hfill Angebot 2 $=$\linie{30mm}
\end{gbteil}
\end{pfaufg}
\begin{pfaufg}{}{}{}
% 2023-OS-K3b-m1 (Original 2023-OS-K3b, Schwere II, rang 15; herausgelöste Fassung: keine)
\begin{gbteil}{b)}{}{FOR}
Herr Mert mietet den Transporter von Montag bis Freitag und fährt insgesamt 450 km. Berechne die Kosten für beide Angebote. Gib an, welches billiger ist. Stelle für Angebot 2 eine Gleichung auf, mit der man die Kosten einer Woche für jede Kilometerzahl berechnen kann.
\karo{3}
\end{gbteil}
\fusshilfe{\gfuss{11}{a)~C; B; b)~170,50 €; 151,00 €; Angebot 2 günstiger; 0,08x + 115}}
\end{pfaufg}

\gmehr{14}
\end{document}